\section{Introduction}

Rubin Observatory is being built to observe the sky repeatedly and rapidly. This creates a
different type of operations problem. The observing system must work as one chain: telescope and
camera, summit and base infrastructure, high-speed networks, storage systems, computing clusters,
databases, control services, and data processing services. If one part of the chain is slow or
becomes unstable, the effect can be seen elsewhere. A storage problem can look like an application
problem. A network issue can look like a data-acquisition issue. A problem in a tiny computing
process can affect a dashboard used by a scientist or an operator during the night.

Hence we made the observability platform part of the observatory itself. We are observing the
system that observes the sky. To conduct reliable scientific operations, we need to view the
observatory with the same seriousness with which we view the sky.

Rubin Observatory will push up to 10\,TB of data through the pipeline every night, triggering
millions of alerts \cite{rubin-dm,rubin-keynumbers,rubin-alerts}. Therefore, decisions have to
happen fast; there is no time to wait when the next image is arriving in 34 seconds. A minute lost
chasing down a graph is a minute the telescope cannot get back.

The observability system described in this paper is not a whiteboard architecture; it reflects
years of operational experience, shaped by real incidents, evolving requirements, and the kind of
hard-won knowledge that only comes from running infrastructure under pressure. The design decisions
made capture not just what was built, but also how the team thinks about keeping environments
consistent and ensuring changes are traceable \cite{lsst-it-k8s-cookbook}.